\begin{python}
    from scripts.protocol import *
\end{python}

\section*{1. Исследование искажений в цепи первого порядка}
\begin{longtblr}[
    caption = {Параметры сигналов в цепи первого порядка},
    label = {tab:first_order_data},
]{
    colspec = { Q[c,m] Q[c,m] Q[c,m] X[l,m] },
    hlines, vlines,
    rowhead = 1,
}
    $t_{\text{и}}$, мкс & $U_{\text{вх max}}$, В & $U_{\text{вых max}}$, В & Визуальная оценка искажений (фронты, спад) \\
    \pv*[-6.0]{p['fod']['ti'][0]} & \pv*[.0]{p['fod']['Uin'][0]} & \pv*[.2]{p['fod']['Uoutmax'][0]} & Сильное сглаживание: импульс не успевает достичь максимума, вершина острая.\\
    \pv*[-6.0]{p['fod']['ti'][1]} & \pv*[.0]{p['fod']['Uin'][1]} & \pv*[.2]{p['fod']['Uoutmax'][1]} & Искажения в виде сглаженных углов. Импульс успевает достицчь установившегося значения, вершина плоская.\\
\end{longtblr}

\section*{2. Исследование искажений в цепи второго порядка}
\begin{figure}[H]
    \centering
    \includegraphics[width=0.75\textwidth]{figs/4000Om.jpg}
    \caption{4000 Ом}
    \label{fig:4000Om.jpg}
\end{figure}
\begin{figure}[H]
    \centering
    \includegraphics[width=0.75\textwidth]{figs/670Om.jpg}
    \caption{670 Ом}
    \label{fig:670Om.jpg}
\end{figure}
\begin{figure}[H]
    \centering
    \includegraphics[width=0.75\textwidth]{figs/100Om.jpg}
    \caption{100 Ом}
    \label{fig:100Om.jpg}
\end{figure}

\section*{3. Исследование искажений в цепи высокого порядка}
\begin{longtblr}[
    caption = {Амплитудно-частотная характеристика цепи высокого порядка},
    label = {tab:bode_data},
]{
    colspec = { X[c,m] X[c,m] },
    hlines, vlines,
    rowhead = 1,
}
    $f$, кГц & $|H_U(j\omega)|$ \\
    \expanded{\pt*[3.0, .2]{
        p['bd']['f'],
        p['bd']['Hu']
    }}
\end{longtblr}

\begin{figure}[H]
    \centering
    \includegraphics[width=0.6\textwidth]{figs/high_order.jpg}
    \caption{Высоких порядоков}
    \label{fig:high_order.jpg}
\end{figure}